
\documentclass{article}
\usepackage{colortbl}
\usepackage{makecell}
\usepackage{multirow}
\usepackage{supertabular}

\begin{document}

\newcounter{utterance}

\centering \large Interaction Transcript for game `bbh', experiment `dyck\_languages', episode 108 with reflexion{-}llama{-}3.1{-}8b{-}instruct.
\vspace{24pt}

{ \footnotesize  \setcounter{utterance}{1}
\setlength{\tabcolsep}{0pt}
\begin{supertabular}{c@{$\;$}|p{.15\linewidth}@{}p{.15\linewidth}p{.15\linewidth}p{.15\linewidth}p{.15\linewidth}p{.15\linewidth}}
    \# & \multicolumn{2}{c}{Player} && \multicolumn{2}{c}{Game Master} \\
    \hline

    \theutterance \stepcounter{utterance}  
    & & & \multicolumn{4}{p{0.6\linewidth}}{
        \cellcolor[rgb]{0.9,0.9,0.9}{
            \makecell[{{p{\linewidth}}}]{
                \texttt{\tiny{[P1$\langle$GM]}}
                \texttt{You have attempted this same task before and did not succeed. These are the plans you wrote after those attempts.} \\
\texttt{Trial \#1: New plan: I failed to correctly close the Dyck{-}n word because I didn't consider the nesting of the parentheses correctly. In the input < \{ < ( ), I started by closing the innermost parentheses (, but then I closed the \{ before the <. This led to an incorrect completion of the sequence. Next time, I will carefully examine the nesting of the parentheses and close the innermost parentheses first, then work my way outwards. I will also pay closer attention to the order of the closing parentheses.} \\
\\ 
\texttt{Trial \#2: New plan: I failed to correctly close the Dyck{-}n word because I didn't consider the nesting of the parentheses correctly. In the input < \{ < ( ), I started by closing the innermost parentheses (, but then I closed the \{ before the <. This led to an incorrect completion of the sequence. Next time, I will carefully examine the nesting of the parentheses and close the innermost parentheses first, then work my way outwards. I will also pay closer attention to the order of the closing parentheses. However, I should also consider the type of brackets and parentheses used in the input. In this case, the input contains both angle brackets < > and curly brackets \{ \}. I should prioritize closing the angle brackets first, as they are the outermost brackets, and then move to the curly brackets.} \\
\\ 
\texttt{Trial \#3: New plan: I failed to correctly close the Dyck{-}n word because I didn't consider the nesting of the parentheses correctly. In the input < \{ < ( ), I started by closing the innermost parentheses (, but then I closed the \{ before the <. This led to an incorrect completion of the sequence. Next time, I will carefully examine the nesting of the parentheses and close the innermost parentheses first, then work my way outwards. I will also pay closer attention to the order of the closing parentheses. However, I should also consider the type of brackets and parentheses used in the input. In this case, the input contains both angle brackets < > and curly brackets \{ \}. I should prioritize closing the angle brackets first, as they are the outermost brackets, and then move to the curly brackets. Additionally, I should pay attention to the nesting of the angle brackets, as they can be nested within each other. For example, in the input < {[} < {[} \{ < {[} {]} < \{ \} > > \} {]} > \{ \{ ( ) \} \{ < {[} < > {]} > \}, I should close the inner angle brackets < {[} < {[} first, before moving to the outer angle brackets < >.} \\
\\ 
\\ 
\texttt{Correctly close a Dyck{-}n word.} \\
\\ 
\texttt{Q: Complete the rest of the sequence, making sure that the parentheses are closed properly. Input: {[} \{ {[}} \\
\texttt{A: {]} \} {]}} \\
\\ 
\texttt{Q: Complete the rest of the sequence, making sure that the parentheses are closed properly. Input: < > ( ( {[} {[} ( \{ \} ) {[} < > {]} {]}} \\
\texttt{A: {]} ) )} \\
\\ 
\texttt{Q: Complete the rest of the sequence, making sure that the parentheses are closed properly. Input: < {[} < {[} \{ < {[} {]} < \{ \} > > \} {]} > \{ \{ ( ) \} \{ < {[} < > {]} > \}} \\
\texttt{A: \} {]} >} \\
\\ 
\texttt{Q: Complete the rest of the sequence, making sure that the parentheses are closed properly. Input: < \{ < ( )} \\
\texttt{A:} \\
            }
        }
    }
     \\ \\

    \theutterance \stepcounter{utterance}  
    & \multicolumn{4}{p{0.6\linewidth}}{
        \cellcolor[rgb]{0.561,0.737,0.561}{
            \makecell[{{p{\linewidth}}}]{
                \texttt{\tiny{[P1$\rangle$GM]}}
                \texttt{To correctly close the Dyck{-}n word, I will follow the plan I wrote in Trial \#3. } \\
\\ 
\texttt{For the input < \{ < ( ), I will carefully examine the nesting of the parentheses and close the innermost parentheses first, then work my way outwards. I will also pay closer attention to the order of the closing parentheses. } \\
\\ 
\texttt{1. The innermost parentheses is (, so I will start by closing it: < \{ < ( ) {-}> < \{ < >.} \\
\texttt{2. Next, I will close the curly brackets \{ \} in the correct order: < \{ < > \} {-}> < \{ \} >.} \\
\texttt{3. Finally, I will close the angle brackets < >: < \{ \} > {-}> < >.} \\
\\ 
\texttt{The correct completion of the sequence is: < >.} \\
            }
        }
    }
    & & \\ \\

    \theutterance \stepcounter{utterance}  
    & & & \multicolumn{2}{p{0.3\linewidth}}{
        \cellcolor[rgb]{0.9,0.9,0.9}{
            \makecell[{{p{\linewidth}}}]{
                \texttt{\tiny{[GM$|$GM]}}
                \texttt{Target: > \} >} \\
            }
        }
    }
    & & \\ \\

    \theutterance \stepcounter{utterance}  
    & & & \multicolumn{2}{p{0.3\linewidth}}{
        \cellcolor[rgb]{0.9,0.9,0.9}{
            \makecell[{{p{\linewidth}}}]{
                \texttt{\tiny{[GM$|$GM]}}
                \texttt{game\_result = LOSE} \\
            }
        }
    }
    & & \\ \\

    \theutterance \stepcounter{utterance}  
    & & \multicolumn{4}{p{0.6\linewidth}}{
        \cellcolor[rgb]{0.95,0.95,0.95}{
            \makecell[{{p{\linewidth}}}]{
                \texttt{\tiny{[P1$\leftrightarrow$P]}}
                \texttt{Post{-}episode Reflexion (not sent to GM):} \\
\texttt{New plan: I failed to correctly close the Dyck{-}n word because I didn't consider the nesting of the parentheses correctly. In the input < \{ < ( ), I started by closing the innermost parentheses (, but then I closed the \{ before the <. This led to an incorrect completion of the sequence. Next time, I will carefully examine the nesting of the parentheses and close the innermost parentheses first, then work my way outwards. I will also pay closer attention to the order of the closing parentheses. However, I should also consider the type of brackets and parentheses used in the input. In this case, the input contains both angle brackets < > and curly brackets \{ \}. I should prioritize closing the angle brackets first, as they are the outermost brackets, and then move to the curly brackets. Additionally, I should pay attention to the nesting of the angle brackets, as they can be nested within each other. } \\
\\ 
\texttt{However, in my previous attempts, I repeated the same strategy with slight modifications, but I didn't consider the fact that the input contains a mix of different types of brackets and parentheses, and that some of them are nested within each other. } \\
\\ 
\texttt{Next time, I will take a different approach. I will start by identifying the different types of brackets and parentheses in the input, and then I will create a mental map of their nesting structure. I will then use this map to guide my closing of the brackets and parentheses, making sure to close the innermost ones first and working my way outwards. } \\
\\ 
\texttt{For example, in the input < \{ < ( ), I will first identify the different types of brackets and parentheses: angle brackets < >, curly brackets \{ \}, and parentheses ( ). I will then create a mental map of their nesting structure: the angle brackets < > are the outermost brackets, the curly brackets \{ \} are nested within the angle brackets, and the parentheses ( ) are nested within the curly brackets. } \\
\\ 
\texttt{Using this map, I will then close the brackets and parentheses in the correct order: first, I will close the innermost parentheses (, then the curly brackets \{ \}, and finally the angle brackets < >. } \\
\\ 
\texttt{By taking this different approach, I hope to avoid the mistakes of my previous attempts and correctly close the Dyck{-}n word.} \\
            }
        }
    }
    & \\ \\

\end{supertabular}
}

\end{document}
